\documentclass[12pt]{article}
\usepackage[margin=1in]{geometry}
\usepackage{enumitem}
\usepackage{titlesec}
\usepackage{graphicx}
\usepackage{amsmath}
\usepackage{amssymb}

\setlength{\parindent}{0pt}
\setlength{\parskip}{0.5em}

\begin{document}

\begin{center}
    {\LARGE \textbf{IcePEAK}}\\[0.3em]
    {\large CIS 5680 --- Game Design Document}\\[0.3em]
    {\normalsize Group Game Project 2 --- Spring 2026}
\end{center}

\vspace{0.5em}
\textbf{Group Members:} Cecilia Chen, Jacky Park, Yiding Tian

\textbf{Game Title:} IcePEAK

\newpage

%% ============================================================
\section{Executive Summary}
%% ============================================================

% Most important page of the design doc.
% Usually 20--25 lines. Include 1--2 sentence descriptions of high concept,
% target audience, market appeal, game play, implementation details, schedule.

\newpage

%% ============================================================
\section{Game Design --- Creative}
%% ============================================================

% -----------------------------------------------------------------
\subsection{High Concept}
% What is the basic concept of your game?
Conquer a treacherous frozen mountain in VR by driving ice picks into glacial walls with your own hands, managing supporting gadgets as you race against avalanches and crumbling ice to reach the summit alive.

% -----------------------------------------------------------------
\subsection{Design Goals}

\subsubsection{Main Design Features}

IcePEAK is a first-person VR climbing game in which the player physically swings ice picks into a frozen mountainside, pulls themselves upward, and manages a limited set of gadgets to survive environmental hazards on the way to the summit.

\paragraph{Player goals and objectives.}

\textbf{Main goal.} Reach the summit of the mountain alive.

\textbf{Secondary goals.} Discover hidden supply caches along alternate routes; complete the ascent using as few gadgets as possible; find the fastest path to the top.

\textbf{Type of challenges.} Physical dexterity (swinging and placing ice picks accurately under time pressure), route-finding (choosing which path upward balances speed, safety, and resource availability), and resource management (deciding when to spend limited gadgets).

\textbf{Type of conflicts.} Player versus environment. The mountain itself is the antagonist: ice surfaces shatter roughly two seconds after an ice pick is planted, falling rocks threaten the climber, ice storms reduce visibility, and rock surfaces reject ice picks entirely.

\textbf{Winning condition.} The player wins by reaching the summit platform at the top of the mountain. The player respawns at nearest rock platform checkpoint when both hands lose support.

\paragraph{Main rules and procedures.}

\textbf{Operational rules.} The player holds one ice pick in each hand via the Quest~3 motion controllers. To climb, the player swings a pick into an ice surface, grips the controller trigger to hold on, then reaches upward with the other hand and repeats. Ice surfaces crack and break approximately two seconds after a pick is driven in, so the player must keep moving. Rock surfaces are permanent and safe to stand on but cannot be gripped by ice picks. Gadgets are grabbed from the player's body-mounted harness and activated with the trigger.

\textbf{Main game mechanic --- second by second.} Swing pick into ice, grip, reach with the other hand, swing again. Watch the crack pattern spreading from the current hold and listen for the audio warning that it is about to shatter. Release before it breaks and latch onto the next surface.

\textbf{Main game mechanic --- minute by minute.} Survey the rock face ahead from a safe platform. Identify the next resting ledge and plan a route through the ice to reach it. Decide whether to use a gadget (e.g.\ cold spray to buy extra seconds on a critical hold, or a grappling hook to skip a dangerous overhang). Begin the climb, react to dynamic hazards (rockfall, wind gusts), and reach the next platform.

\textbf{What makes this fun.} The two-second ice timer turns every handhold into a micro-deadline, creating a rhythm of tension and relief. Because the player's real arms do the climbing, success feels physically earned. Gadget decisions add a strategic layer: using a gadget now makes life easier but leaves fewer options for harder terrain above.

\paragraph{Player resources.}

The player carries a limited inventory of gadgets, each with a single or limited use:

\begin{itemize}[leftmargin=2em]
    \item \textbf{Grappling gun hook} --- fires a hook that attaches to a distant surface, letting the player zip to a new position. Single use per hook.
    \item \textbf{Cold spray} --- sprayed onto an ice surface to make it more rigid, significantly extending the time before it cracks (e.g.\ from two seconds to six seconds). Single use per canister.
    \item \textbf{Helmet} --- absorbs one falling-rock strike that would otherwise injure or dislodge the player. Breaks after one hit.
    \item \textbf{Poles} --- metal stakes that can be driven into rock surfaces, creating an anchor point that an ice pick can hook onto. This lets the player grip rock walls that would otherwise be unclimbable.
\end{itemize}

Gadgets are found in supply caches scattered across the mountain, often along riskier or longer routes. The player uses them to overcome obstacles that would be impossible or extremely difficult with ice picks alone.

\paragraph{Boundaries and constraints.}

The game takes place on a single massive frozen mountain. The player can only move by physically climbing; there is no teleportation or artificial locomotion. Ice surfaces are temporary (they shatter after being struck), and rock surfaces are permanent but cannot be gripped by picks unless a pole is attached. The player cannot carry unlimited gadgets --- the harness has a fixed number of slots. Falling below the last safe rock platform results in death and a restart from that platform. Visibility, wind, and ice fragility worsen with altitude, naturally constraining the pace of the ascent.

\subsubsection{Appeal}

\textbf{Genre.} First-person VR survival-adventure with physics-based climbing and resource management.

\textbf{Target audience.} VR gamers aged 16--35 who enjoy physically active experiences and survival/adventure games. Players of titles like \textit{The Climb} (Crytek), \textit{Until You Fall}, and \textit{The Walking Dead: Saints \& Sinners} who are looking for a more intense, physics-driven challenge.

\textbf{Why it is fun.} No existing VR climbing game forces the player to race against crumbling holds. The destructible-ice mechanic transforms climbing from a relaxed puzzle into a high-stakes rhythm game played with the whole body. The gadget system adds strategic depth without the abstraction of stamina bars or health meters --- every tool is a tangible object the player grabs off their own harness and physically uses. The result is a game that is immediately intuitive (swing, grab, climb) yet deeply replayable (route optimization, gadget conservation, speed runs).

\subsubsection{Look and Feel}

\textbf{Art style.} Semi-realistic with stylized lighting. The mountain features stark contrasts between blue-white ice, dark grey rock, and warm amber light from supply caches and the distant sunrise above the summit. The palette shifts from cool blues at the base to harsh whites and deep purples near the peak to convey increasing altitude and danger.

\textbf{Color scheme.} Dominant cool tones (ice blue, slate grey, snow white) punctuated by warm accent colors on gadgets and supply caches (orange, amber) so interactive objects are immediately readable.

\textbf{Similar games.} Visually inspired by \textit{The Climb} and \textit{Steep}. The tension and survival feel draw from \textit{Celeste} (translated into 3D/VR) and the environmental storytelling of \textit{Journey}.

% -----------------------------------------------------------------
\subsection{Worlds, Characters and Story}

\subsubsection{Back Story}

There is no traditional narrative cutscene. The player wakes at a base-camp ledge partway up the mountain with ice picks already in hand and a loaded harness. Scattered environmental details --- a torn flag, an abandoned tent half-buried in snow, scratched tally marks on a rock wall --- hint that others have attempted this climb before and failed. The story is the ascent itself: why the climber is here is left to the player's imagination; the only thing that matters is reaching the top.

This minimal backstory is revealed passively through the environment rather than dialogue or text, keeping the player focused on the physical act of climbing.

\subsubsection{Spaces/Worlds}

The game takes place on a single massive frozen mountain divided into three altitude zones, each with a distinct visual identity and escalating difficulty:

\begin{itemize}[leftmargin=2em]
    \item \textbf{Zone 1 --- Lower Ice Field (base camp to ~30\% elevation).} Wide ice walls with generous spacing between rock platforms. Ice is thick and cracks slowly (closer to three seconds). Falling rocks are infrequent. Supply caches are common. This zone serves as the tutorial area where the player learns the core climbing loop and gadget usage.

    \item \textbf{Zone 2 --- The Ridge (~30\%--70\% elevation).} Narrower climbing corridors, steeper overhangs, and longer gaps between safe platforms. Ice cracks at the standard two-second rate. Rockfall events become regular, and brief ice storms temporarily reduce visibility. Alternate routes diverge here, offering meaningful risk--reward choices.

    \item \textbf{Zone 3 --- The Summit Push (~70\%--100\% elevation).} Near-vertical faces, extremely fragile ice (under two seconds), persistent wind that sways the player's view, and reduced visibility from blowing snow. Supply caches are rare. This zone demands precise gadget management and fast, confident climbing.
\end{itemize}

% TODO: Include sketches of each zone.

\subsubsection{Characters}

There is one character: the player. IcePEAK is a solo first-person experience with no NPCs. The player sees their own hands gripping ice picks, a harness across their chest with gadget slots, and climbing gear on their body. The character has no spoken dialogue; the player's physical actions \emph{are} the character.

% TODO: Include sketch of the player's first-person body view (hands, harness, gear).

\subsubsection{Levels of Difficulty}

Difficulty increases naturally with altitude rather than through a menu setting:

\begin{itemize}[leftmargin=2em]
    \item \textbf{Ice fragility} decreases with elevation --- holds last shorter before shattering.
    \item \textbf{Distance between rock platforms} increases, requiring longer uninterrupted climbing stretches.
    \item \textbf{Environmental hazards} become more frequent and intense (more rockfall, stronger wind, ice storms).
    \item \textbf{Supply caches} become sparser, forcing the player to be more conservative with gadgets.
    \item \textbf{Route complexity} increases --- paths branch more and include dead ends, overhangs, and traversals that require gadgets to pass.
\end{itemize}

This elevation-driven curve means the player never encounters an artificial difficulty screen; the mountain itself gets harder the higher they go.

% -----------------------------------------------------------------
\subsection{Interaction Models}

\subsubsection{User Interface --- Navigation and Movement}

There is no traditional HUD or menu overlay during gameplay. All information is communicated through the player's body and the environment:

\begin{itemize}[leftmargin=2em]
    \item \textbf{Navigation.} The player moves exclusively by physically climbing --- swinging ice picks into surfaces and pulling themselves up. There is no joystick locomotion, teleportation, or artificial movement. On rock platforms the player can walk freely within the platform bounds using room-scale tracking.
    \item \textbf{Gadget selection.} Gadgets are stored on the player's body harness. The player looks down at their chest/belt, reaches toward a gadget slot, and grabs the item with the grip button. No radial menus or button-based inventory.
    \item \textbf{Pause menu.} Accessed by pressing the Quest~3 menu button. Pauses the game and displays options (resume, restart from last checkpoint, quit).
\end{itemize}

% TODO: Include mockup sketch of first-person harness view with gadget slots labeled.

\subsubsection{Game Play Sequence and Levels}

The game follows a linear upward progression through three zones:

\begin{enumerate}[leftmargin=2em]
    \item \textbf{Base Camp Start.} Player spawns on a rock platform. Brief environmental cues teach the controls: an ice wall directly ahead with a rock ledge just above it invites the first climb. A supply cache with one of each gadget is nearby.
    \item \textbf{Zone 1 --- Lower Ice Field.} Gentle climbing with wide holds and frequent platforms. The player learns the ice-cracking mechanic and practices using gadgets in low-stakes situations.
    \item \textbf{Zone 2 --- The Ridge.} Routes branch. The player chooses between paths with different risk--reward profiles. Hazards (rockfall, ice storms) are introduced. Longer stretches between checkpoints.
    \item \textbf{Zone 3 --- The Summit Push.} Intense vertical climbing with fragile ice, rare supplies, and persistent environmental threats. Culminates in a final exposed face leading to the summit.
    \item \textbf{Summit.} The player pulls themselves onto the peak. A panoramic vista rewards the climb. The game ends.
\end{enumerate}

Transitions between zones are seamless --- there are no loading screens. The player simply climbs from one zone into the next, with the environment gradually changing around them.

% TODO: Include flowchart diagram of game sequence.

\subsubsection{User/Environment --- Obstacles and Props}

\textbf{Ice surfaces.} The primary interactive surface. When the player swings an ice pick into ice with sufficient velocity, the pick embeds and the player can hold on. A crack pattern begins spreading from the impact point (visual effect plus audio cue). After approximately two seconds the surface shatters and the hold is lost. Collision detection between the ice pick tip and ice surfaces uses Unity's Rigidbody collision system. Ice destruction is handled by replacing the intact mesh with a pre-fractured version using a shattering script triggered by a timer.

\textbf{Rock surfaces.} Permanent, unbreakable. Ice picks bounce off rock (detected via Rigidbody collision and surface material tag). The player can stand and walk on rock platforms. Poles can be driven into rock to create anchor points.

\textbf{Supply caches.} Glowing containers embedded in the cliff face or resting on ledges. The player grabs them and the contents are automatically added to empty harness slots. Indicated by a warm amber glow and a subtle audio hum to aid discovery.

\textbf{Falling rocks.} Triggered by proximity or timer events. Rocks fall from above along semi-random paths. A shadow and rumbling audio cue warn the player 1--2 seconds before impact. If a rock hits the player without a helmet, the player is knocked off the wall (losing both grips). With a helmet equipped, the first strike is absorbed and the helmet breaks.

\textbf{Ice storms.} Periodic events that reduce visibility (fog/particle effects) and add a slight lateral force to the player's view, simulating wind. During storms the player must climb more carefully as visual cues for cracking ice are harder to see.

\textbf{Wind gusts.} Brief directional forces that nudge the player's camera and add sway, increasing disorientation at higher altitudes.

\textbf{Physics.} All climbing interactions use Unity's physics engine. Ice pick embedding is determined by the velocity of the controller swing (a minimum threshold must be met). The player's vertical position is driven by the physical act of pulling down on an embedded pick (inverse kinematics maps controller position to body movement). Falling uses standard gravity with a brief ragdoll before respawn.

\subsubsection{User/Character}

The player \emph{is} the character. There is no avatar to control indirectly. The Quest~3 controllers map directly to the player's hands:

\begin{itemize}[leftmargin=2em]
    \item \textbf{Grip button} --- hold to maintain grip on an embedded ice pick or a grabbed gadget.
    \item \textbf{Trigger} --- activate a held gadget (fire grappling hook, spray cold spray, place pole).
    \item \textbf{Swing motion} --- the speed and angle of the controller swing determine whether the ice pick embeds successfully and where.
\end{itemize}

The player looks down to see their torso, harness with gadget slots, and hands holding ice picks. Head tracking provides full look-around. There are no other characters to interact with.

\subsubsection{Character/Character}

Not applicable. IcePEAK is a single-player game with no NPCs or AI characters. All conflict is player versus environment.

\subsubsection{Puzzle Design}

IcePEAK does not feature traditional puzzles. The ``puzzle'' element comes from route-finding: the player must read the mountain face and determine a viable path from one rock platform to the next, choosing which ice surfaces to use, which gadgets to spend, and how to sequence their moves before holds shatter. Specific route-puzzle scenarios include:

\begin{itemize}[leftmargin=2em]
    \item \textbf{Overhang gaps} where the player must use a grappling hook or find an alternate traverse.
    \item \textbf{Large rock faces} with no ice, requiring poles to create anchor points before crossing.
    \item \textbf{Timed sequences} where multiple ice surfaces must be crossed in rapid succession before any of them shatter, demanding the player plan the order and direction of their swings in advance.
\end{itemize}

% TODO: Include sketches of route-puzzle scenarios in Section 3.2 Storyboards.

\subsubsection{Motion Tracking}

IcePEAK relies heavily on the Quest~3's 6DoF (six degrees of freedom) tracking for both the headset and controllers:

\begin{itemize}[leftmargin=2em]
    \item \textbf{Controller position and velocity} are used to determine swing direction, speed, and impact point when the player drives an ice pick into a surface. A minimum velocity threshold prevents the player from gently placing picks without swinging.
    \item \textbf{Controller orientation} determines the angle of the ice pick at impact, affecting whether the pick embeds cleanly or glances off.
    \item \textbf{Headset tracking} provides the player's viewpoint and is used for look-based environmental awareness (spotting routes, seeing crack patterns, noticing falling-rock shadows).
    \item \textbf{Grip state} is tracked via the controller's grip button to determine whether the player is holding onto an embedded pick or has released.
\end{itemize}

\subsubsection{Multi-Player}

Not applicable. IcePEAK is a single-player experience.

\subsubsection{Mobile}

Not applicable.

\subsubsection{Networked Play}

Not applicable.

% -----------------------------------------------------------------
\subsection{Performance and Scoring}

\subsubsection{State Variables}

The game consists of a single continuous level with multiple checkpoint platforms. The following variables are necessary to save/restore or pause/resume the session:

\begin{itemize}[leftmargin=2em]
    \item \textbf{Current checkpoint} --- the highest rock-platform checkpoint the player has reached. Used as the respawn point.
    \item \textbf{Current altitude} --- the player's real-time vertical position on the mountain (meters above base camp).
    \item \textbf{Highest altitude} --- the maximum altitude the player has reached during this session.
    \item \textbf{Gadget inventory} --- which gadgets occupy each harness slot (type and count per slot).
    \item \textbf{Helmet status} --- whether the helmet is equipped and intact or already broken.
    \item \textbf{Ice surface states} --- for each ice surface near the player: intact, cracking (with remaining time), or shattered. Surfaces far below the current checkpoint can be reset.
    \item \textbf{Supply cache states} --- which caches have been collected (so they do not respawn).
    \item \textbf{Active hazard timers} --- any in-progress rockfall or ice storm events and their remaining duration.
    \item \textbf{Player hand states} --- whether each hand currently has a pick embedded in a surface (left grip, right grip, or free).
\end{itemize}

On pause, all physics and timers freeze. On resume, they continue from the saved state. On respawn (both hands lose support), the player is placed on the highest reached checkpoint with ice surfaces in that area reset to intact and the gadget inventory preserved as-is.

\subsubsection{Feedback}

\textbf{Positive feedback:}
\begin{itemize}[leftmargin=2em]
    \item \textbf{Successful pick embed} --- a satisfying \emph{thunk} sound and haptic pulse in the controller confirm the hold is secure.
    \item \textbf{Reaching a checkpoint} --- a brief chime and a warm glow emanates from the rock platform, signaling safety. The current altitude is displayed momentarily as carved markings on the rock.
    \item \textbf{Altitude milestone} --- when the player reaches a new highest altitude, a subtle wind-chime audio cue plays, reinforcing upward progress.
    \item \textbf{Gadget use} --- each gadget has a distinct audio-visual payoff (grappling hook zip, cold spray crystallization effect, pole locking into rock with a metallic clang).
\end{itemize}

\textbf{Negative feedback:}
\begin{itemize}[leftmargin=2em]
    \item \textbf{Ice cracking} --- escalating crackle audio and visible fracture lines warn the player the hold is about to give way. Controller haptics intensify as time runs out.
    \item \textbf{Falling rock warning} --- a shadow passes over the player and a deep rumble plays 1--2 seconds before impact. If hit without a helmet, the player is knocked off the wall.
    \item \textbf{Fall and respawn} --- when both hands lose support, the player falls briefly (with vertigo-inducing visual blur) before fading to black and respawning at the last checkpoint. This short fall sequence reinforces the stakes without being punishing --- gadget inventory is preserved and the checkpoint is nearby.
    \item \textbf{Glancing pick strike} --- if the swing is too slow or the angle is wrong, the pick sparks off the surface with a metallic scrape and no embed occurs, giving immediate feedback to adjust technique.
\end{itemize}

\subsubsection{Performance and Progress Metrics}

\textbf{Progress tracking.} The player's performance is measured by two altitude values:
\begin{itemize}[leftmargin=2em]
    \item \textbf{Current altitude} --- the player's present height on the mountain, visible as environmental markings on checkpoint platforms.
    \item \textbf{Highest altitude} --- the peak altitude reached during the session, tracked persistently.
\end{itemize}

\textbf{Winning.} The player wins by reaching the summit platform at the top of the mountain. Upon arrival, the panoramic vista and a final altitude reading confirm completion.

\textbf{Losing.} There is no permanent game-over state. When both hands lose support (both picks dislodge or ice shatters under both holds simultaneously), the player falls and respawns at the highest checkpoint reached. The player keeps all remaining gadgets. Progress is never lost beyond the last checkpoint, but supply caches do not respawn, so reckless play gradually depletes available resources for the climb ahead.

\newpage

%% ============================================================
\section{Game Design --- Implementation Details}
%% ============================================================

% -----------------------------------------------------------------
\subsection{Design Assumptions}

\subsubsection{Hardware}

\begin{itemize}[leftmargin=2em]
    \item \textbf{Headset:} Meta Quest~3 (standalone mode). The Snapdragon XR2 Gen~2 processor, 8\,GB RAM, and Adreno 740 GPU are the target hardware baseline.
    \item \textbf{Controllers:} Quest~3 Touch Plus controllers (two, one per hand). 6DoF tracking is required for swing detection and spatial positioning.
    \item \textbf{Play space:} Standing or room-scale VR. A minimum 2\,m $\times$ 2\,m play area is recommended so the player can extend both arms overhead and to the sides while climbing.
\end{itemize}

\subsubsection{Software}

\begin{itemize}[leftmargin=2em]
    \item \textbf{Game engine:} Unity 2022 LTS (or newer 2022.3.x patch).
    \item \textbf{XR framework:} Meta XR SDK (formerly Oculus Integration) for Quest~3 input, hand tracking fallback, and performance profiling.
    \item \textbf{Rendering:} Universal Render Pipeline (URP) configured for mobile VR (single-pass instanced rendering, foveated rendering enabled).
    \item \textbf{Target platform:} Android (Quest~3 standalone APK). No PC tethering required.
    \item \textbf{Minimum OS:} Quest system software v62 or later.
\end{itemize}

\subsubsection{Algorithms and Techniques}

\begin{itemize}[leftmargin=2em]
    \item \textbf{Ice pick collision detection.} Unity Rigidbody physics with trigger colliders on the ice pick tips. On collision, the script reads the controller velocity from \texttt{OVRInput} (or the XR Interaction Toolkit velocity provider). If the velocity magnitude exceeds a tunable threshold and the surface tag is \texttt{Ice}, the pick embeds; if the tag is \texttt{Rock}, the pick bounces off with a spark particle effect.

    \item \textbf{Ice destruction system.} Each ice surface is a prefab with an intact mesh and a pre-fractured variant (created with a mesh-fracture tool such as RayFire or a custom Voronoi shattering script). When a pick embeds, a coroutine starts a timer (default 2\,s). At expiry the intact mesh is swapped for the fractured pieces, which are given small outward Rigidbody forces and destroyed after a short lifetime to reclaim memory.

    \item \textbf{Cold spray extension.} Applying cold spray to a surface multiplies the crack timer (e.g.\ $2\text{s} \rightarrow 6\text{s}$). Implemented by overriding the timer duration on that specific ice prefab instance and playing a frost particle effect.

    \item \textbf{Grappling hook.} On trigger press a raycast is fired from the hook prop in the direction the player is aiming. If it hits a valid surface within range, a line renderer draws the rope and the player's position is lerped toward the hit point over $\sim$0.5\,s. At arrival the hook is consumed.

    \item \textbf{Pole placement.} When the player holds a pole against a rock surface and presses trigger, a raycast confirms the surface tag is \texttt{Rock}. The pole prefab is instantiated at the hit point with a fixed joint to the rock. The pole's tip has an \texttt{Ice}-tagged collider so that ice picks can embed in it.

    \item \textbf{Player body / inverse kinematics.} A simple upper-body IK rig (Final IK plugin or Unity's built-in Animation Rigging package) maps headset and controller transforms to a visible torso, arms, and hands. This provides the harness and gear visibility without full-body tracking.

    \item \textbf{Climbing locomotion.} When a pick is embedded and the player pulls the controller downward, the player's camera rig moves upward by the inverse of the controller displacement. This is handled per-frame by reading the delta of the gripping controller's position and applying the negated vector to the XR Origin transform.

    \item \textbf{Environmental hazards.}
    \begin{itemize}
        \item \emph{Falling rocks:} Spawned by trigger zones or timers. Each rock is a Rigidbody with a collider that follows a semi-random downward path. A shadow projector and audio source are attached and activated 1--2\,s before the rock enters the player's zone.
        \item \emph{Ice storms:} A full-screen particle system and post-processing fog increase. A global wind vector is applied as a small force to the camera rig, simulating buffeting.
    \end{itemize}

    \item \textbf{Checkpoint system.} Each rock platform has a trigger collider. When the player's feet (or camera rig Y-position) enters the trigger, the platform is recorded as the current checkpoint and the highest-altitude value is updated if applicable.

    \item \textbf{Performance optimization.} Level-of-detail (LOD) groups on distant mountain geometry, shader-based fog to limit draw distance, dynamic batching of ice fragment meshes, and object pooling for rock hazards and particle effects to maintain 72\,fps on Quest~3.
\end{itemize}

% -----------------------------------------------------------------
\subsection{Storyboards}

% TODO: Insert storyboard sketches here. The following descriptions indicate
% what each storyboard frame should depict.

\textbf{Frame 1 --- Base Camp.} The player's first-person view from the starting rock platform. Ice picks are visible in both hands. The harness with gadget slots is visible on the chest. A short ice wall leads to the first ledge above.

\textbf{Frame 2 --- First Climb.} The player mid-swing, driving a pick into ice. Crack lines are beginning to spread from a previous hold below. The next rock ledge is two body-lengths above.

\textbf{Frame 3 --- Gadget Use (Cold Spray).} The player holds a cold-spray canister against an ice surface. A frost effect radiates from the spray point. The ice appears thicker and bluer, indicating the extended hold timer.

\textbf{Frame 4 --- Rockfall Hazard.} A shadow falls across the player's hands. Above, a boulder is dislodging from the cliff. The player must decide whether to swing laterally to avoid it or trust the helmet.

\textbf{Frame 5 --- Route Choice (Zone 2).} From a wide rock platform, two paths are visible: a steep direct route with sparse holds on the left, and a longer traverse to the right that passes a glowing supply cache. Both paths converge at a ledge above.

\textbf{Frame 6 --- Summit Push (Zone 3).} Near-vertical ice face, blowing snow particles obscuring the view. The player's hands grip picks embedded in thin, heavily cracked ice. The summit ridge is faintly visible above through the storm.

\textbf{Frame 7 --- Summit.} The player pulls onto the peak. A panoramic mountain vista stretches in every direction. The sky is clear. Altitude markings on a cairn confirm the summit has been reached.

% -----------------------------------------------------------------
\subsection{Design Logic}

\subsubsection{Finite State Machines --- Events/States/Effects}

% TODO: Include FSM diagram. The following describes the states and transitions.

\textbf{Player FSM:}

\begin{itemize}[leftmargin=2em]
    \item \textbf{Idle (on platform)} --- Player is standing on a rock platform. Can look around, access harness, grab gadgets. \emph{Transition:} pick embeds in surface $\rightarrow$ Climbing.
    \item \textbf{Climbing} --- One or both picks are embedded. Player is actively ascending. Ice timers are running. \emph{Transitions:} both picks free + on platform $\rightarrow$ Idle; both picks free + not on platform $\rightarrow$ Falling; checkpoint trigger entered $\rightarrow$ Idle (on platform).
    \item \textbf{Using Gadget} --- Player has activated a gadget (brief state). \emph{Transitions:} grappling hook $\rightarrow$ Zipping; cold spray / pole $\rightarrow$ return to Climbing or Idle.
    \item \textbf{Zipping} --- Player is being pulled toward grappling hook target. Movement is automated. \emph{Transition:} arrival at target $\rightarrow$ Climbing or Idle (depending on surface type).
    \item \textbf{Falling} --- Both hands lost support and player is not on a platform. Gravity takes over. \emph{Transition:} after brief fall $\rightarrow$ Respawning.
    \item \textbf{Respawning} --- Screen fades to black, player is repositioned at last checkpoint. Nearby ice surfaces are reset to intact. \emph{Transition:} fade-in complete $\rightarrow$ Idle (on platform).
\end{itemize}

\textbf{Ice Surface FSM:}

\begin{itemize}[leftmargin=2em]
    \item \textbf{Intact} --- Default state. Surface can receive a pick embed. \emph{Transition:} pick embeds $\rightarrow$ Cracking.
    \item \textbf{Cracking} --- Timer is running ($\sim$2\,s default). Crack lines spread visually; audio escalates. \emph{Transitions:} timer expires $\rightarrow$ Shattered; cold spray applied $\rightarrow$ Cracking (timer extended).
    \item \textbf{Shattered} --- Surface breaks into fragments that fall away. No longer usable. \emph{Transition:} player respawns at checkpoint $\rightarrow$ Intact (for surfaces near the checkpoint only).
\end{itemize}

\textbf{Hazard FSM (Falling Rock):}

\begin{itemize}[leftmargin=2em]
    \item \textbf{Dormant} --- Rock is part of the environment, not yet triggered. \emph{Transition:} player enters trigger zone or timer fires $\rightarrow$ Warning.
    \item \textbf{Warning} --- Shadow projector activates, rumble audio plays (1--2\,s). \emph{Transition:} warning time elapsed $\rightarrow$ Falling.
    \item \textbf{Falling} --- Rock Rigidbody is enabled, rock falls under gravity. \emph{Transitions:} hits player $\rightarrow$ Impact (check helmet); exits play area $\rightarrow$ Destroyed.
    \item \textbf{Impact} --- If helmet equipped: helmet breaks, player stays attached. If no helmet: player enters Falling state (both grips released). Rock is destroyed.
\end{itemize}

% TODO: Include FSM graph diagrams (hand-drawn or Visio).

\subsubsection{User Solution/Actions}

The player wins by climbing from base camp to the summit, passing through all three altitude zones. The core loop at every stage is: survey the route from a safe platform, plan the sequence of holds and gadget usage, then execute the climb.

\textbf{Key challenge scenarios and solutions:}

\begin{enumerate}[leftmargin=2em]
    \item \textbf{Standard ice traverse.} The player swings from hold to hold, staying ahead of the 2-second crack timer. Solution: maintain a steady rhythm of swing--grip--reach--swing, always placing the next pick before the current hold shatters.

    \item \textbf{Long gap between platforms.} A stretch of ice too long to cross before early holds shatter. Solution: use cold spray on one or two critical holds midway through to buy extra time, creating ``rest points'' within the ice.

    \item \textbf{Rock wall barrier.} A section of pure rock with no ice. Ice picks cannot grip. Solution: drive a pole into the rock to create an anchor, then hook the ice pick onto the pole to traverse across.

    \item \textbf{Overhang or distant ledge.} A gap too wide to reach by climbing. Solution: fire the grappling hook at a surface across the gap and zip over.

    \item \textbf{Rockfall gauntlet.} A zone where rocks fall frequently. Solution: equip the helmet for insurance, watch for shadow cues, and time climbing bursts between rockfall waves. Alternatively, find a sheltered route (e.g.\ under an overhang) that avoids the rockfall zone.

    \item \textbf{Ice storm section.} Reduced visibility makes it hard to spot the next hold. Solution: move cautiously, listen for audio cues (crack sounds indicate nearby ice surfaces), and rely on the warm glow of supply caches as waypoints.
\end{enumerate}

The overall strategy is to conserve gadgets for the harder upper zones while using the lower zones to learn the climbing mechanics and build confidence.

% -----------------------------------------------------------------
\subsection{Software Versions}

\subsubsection{Alpha Version Features (vertical slice through total experience)}

The alpha is the first playable build, targeting the April~20 playtest. It delivers a vertical slice: a short but complete climb from a starting platform to a summit marker, demonstrating the core loop end-to-end.

\begin{itemize}[leftmargin=2em]
    \item Core climbing mechanic: swing ice picks into ice, grip, pull up. Physics-based embed detection with velocity threshold.
    \item Ice destruction: 2-second crack timer, visual crack effect, surface shattering into fragments.
    \item Rock platforms as checkpoints with respawn on fall.
    \item At least one gadget functional (cold spray or grappling hook).
    \item Basic environment: one simplified mountain column with 3--4 checkpoint platforms and connecting ice surfaces.
    \item Altitude tracking: current and highest altitude displayed at checkpoints.
    \item Haptic feedback on pick embed and ice cracking.
    \item Runs on Quest~3 standalone at 72\,fps.
\end{itemize}

\subsubsection{Beta Version Features}

The beta targets the April~27 playtest. It expands the alpha into the full game experience.

\begin{itemize}[leftmargin=2em]
    \item Full mountain environment with three zones (Lower Ice Field, The Ridge, The Summit Push) and all checkpoint platforms.
    \item All four gadgets functional: grappling hook, cold spray, helmet, poles.
    \item Body-mounted harness with physical gadget grab from slots.
    \item Supply caches placed throughout the mountain.
    \item Route branching in Zone~2 with alternate paths.
    \item Environmental hazards: falling rocks with shadow/audio warning, ice storms with particle effects and wind force.
    \item Difficulty scaling with altitude (ice fragility, hazard frequency, supply scarcity).
    \item Audio pass: ambient wind, crack escalation sounds, gadget sound effects, rockfall rumble.
    \item Summit finale with panoramic vista.
    \item Performance optimization pass (LODs, object pooling, foveated rendering).
\end{itemize}

\subsubsection{Descriptions of any tutorial levels and/or self-running Demos}

There is no separate tutorial level. The base-camp area at the start of the game serves as an implicit tutorial:

\begin{itemize}[leftmargin=2em]
    \item The first ice wall is short (one body-length) with a rock ledge directly above, teaching the basic swing--grip--pull cycle with no risk.
    \item A supply cache at base camp contains one of each gadget. Nearby surfaces are arranged so the player can experiment --- a rock wall for pole practice, a distant ledge for grappling hook use, and a thick ice surface for cold spray.
    \item Ice in Zone~1 cracks slightly slower ($\sim$3\,s) to give new players extra reaction time before the standard 2-second timer kicks in at Zone~2.
\end{itemize}

No self-running demo is planned. The gameplay trailer (due May~8) will serve as the promotional demonstration.

\newpage

%% ============================================================
\section{Work Plan}
%% ============================================================

% -----------------------------------------------------------------
\subsection{Tasks}
% List and number ALL the tasks and subtasks which are necessary to develop
% your game application.
% Provide separate descriptions of each task and subtask, which members of
% the group are assigned to it, and the expected task duration (in days).

% -----------------------------------------------------------------
\subsection{Milestones}

\subsubsection{Minor}
% On a week by week basis describe the functionality you plan to achieve
% (and will be able to demonstrate).

\subsubsection{Major}
% Describe the major milestones of the project.

\paragraph{Alpha Version.}
% Describe the main functionality you plan to achieve and demonstrate.

\paragraph{Beta Version.}
% Describe the main functionality you plan to achieve and demonstrate.

% -----------------------------------------------------------------
\subsection{Development Schedule}
% Organize your work plan tasks in the form of a Gantt chart.
% List all the major tasks and subtasks (including durations) and the dates
% of each minor (weekly) milestone and major milestone (e.g.\ alpha, beta, final).

\end{document}
